% PARKED CELL (moved from reviewer2.tex on 2026-10-25).
% No referee point stands behind this text; the \refcom placeholder has been
% removed until real referee wording exists. Reactivate by moving back into
% reviewer2.tex under the next free R2 number and restoring a quote line.
% ---------------------------------------------------------------------
% AG (Oct 2026): PARKED -- this cell has no referee comment behind it. It holds a correction we found
% ourselves while revising: our pricing tables had compared nominal magnitudes in two different units (the
% wage-pinned rows in wage units, the flexible-wage rows in consumption units), and the closure we had called
% `fixed real wage' does not peg the real wage at all -- it pins the money wage and lets employment adjust,
% which is why it is now called employment adjustment. Neither report mentions units, numeraire or
% normalisation, so the quote line cannot be filled. The draft is kept below, commented out, until we decide
% where it belongs; placement suggestions are in the session report.
% \subsection*{R2.13 --- Units of the reported price changes (parked: no referee point behind it)}
% \response{We thank the referee: this was wrong, and chasing it turned up a second,
% more substantive error of ours. The two are related but distinct, and we have
% treated them as such.
%
% \emph{Units.} The pricing tables compared nominal magnitudes expressed in two
% different units. The flexible-wage closures carried the consumption price index as
% numeraire, so their rows were already in consumption units; the wage-pinned rows
% fixed the money wage instead, which left prices in wage units. In sharp contrast, the revised manuscript reports consumption units for all models: prices are divided by the cell's own CPI, the GDP deflator is rescaled with them, and the maximum sectoral deviation is computed on the normalised price vector. 
%
% It remains to be added that the CPI sits at one as long as no prices move and, hence, plays a role mainly for the no-reallocation and the elastic labor supply closure as well as the extensions presented in Section~6.
%
% % AG (Oct 2026): check critically what we really need from the paragraph below.
% \emph{Closure.} The reason the CPI moved at all is that the closure we had called
% a ``fixed real wage'' closure does not peg the real wage. It fixes the money wage
% as the numeraire, sets the net external transfer to zero and drops the
% labour-supply equation, so that employment adjusts. With one production factor
% and constant returns the real wage is not an instrument at all: it is fixed by
% technology, $w/\mathcal{P}(p) = 1/\mathcal{P}(\pi(A^{H}))$, so a peg on it is
% redundant at the calibration point and infeasible away from it (the
% no-real-wage-peg corollary in Appendix~B). We have renamed the closure
% \emph{employment adjustment}, stated its equations exactly as they are solved,
% and added the identity with its proof. No estimate in the paper changes: the
% kernel never solved the peg, and real outcomes are invariant to the nominal
% scale. What changes is the label, the interpretation --- in particular we no
% longer claim that workers' purchasing power is held fixed there --- and the
% units.
%
% The two corrections cannot substitute for one another, which is why we report
% both: renormalising cannot restore a constant real wage, and relabelling cannot
% repair mixed units. Reporting the wage-pinned rows in consumption units is in
% fact what makes the movement of the real wage visible instead of hidden inside
% the reporting convention.}
% \loc{Sec.~4 (the numeraire of each closure); Secs.~5--6 (the affected tables);
% Appendix~A (the nesting claims, scoped to the demand-only slice); Appendix~B
% (the no-real-wage-peg corollary and the accompanying remark).}
